\label{sec:introduction}

When the Ohio Redistricting Commission drew state legislative maps in 2021,
residents of the Cuyahoga River valley testified that their community identity
was shaped more by the physical landscape---the valley floor, the bluffs, the
river itself---than by county or municipal boundaries.
The Commission noted the testimony and then, lacking any computational mechanism
to encode physical land use as a community signal, drew lines based on population
equality and VRA compliance.
The valley communities were divided across three districts that reflected none of
the physical geography that defined them.

This paper provides that computational mechanism.

Physical land use differs from administrative community signals (M.6 zone
co-membership) and economic community signals (M.1 economic character) in an
important respect: it is what actually exists on the ground.
Zoning maps reflect what a municipality \emph{permits}; land use maps reflect
what has actually been built.
A census tract may be zoned for mixed use but be entirely developed as a
residential neighborhood; the NLCD classification captures the residential
reality, not the zoning aspiration.
For redistricting purposes---where the question is what kind of community
actually lives within a proposed district---the physical reality is the
appropriate input.

The M-track community character framework, introduced in M.0
\citep{m0framework}, operationalizes diverse community signals as edge
weights in the census-tract adjacency graph.
Communities that belong together receive heavy edges that METIS avoids
cutting; communities that are naturally distinct receive light edges that
invite cuts at community transitions.
M.2 is the land use implementation paper.
It applies the M-track framework using the USGS National Land Cover Database
(NLCD) 2021, the most comprehensive and authoritative land use classification
of the contiguous United States.

Physical land use is the right M-track signal for several reasons.

\begin{enumerate}
  \item \textbf{Empirical grounding.}
    NLCD is a remote sensing product derived from Landsat satellite imagery.
    It measures what is actually on the ground, not what jurisdictions say
    should be there.
    Classification accuracy for the developed land classes is above 85\%
    based on the MRLC validation protocol \citep{homer2020nlcd,dewitz2023nlcd2021}.
  \item \textbf{Legal recognition.}
    Physical geography---including natural boundaries such as rivers, mountain
    ridges, and coastal features---is one of the oldest recognized redistricting
    criteria.
    Courts have consistently upheld plans that follow natural boundaries
    and struck down plans that cross natural boundaries without geographic
    justification.
  \item \textbf{Water boundary doctrine.}
    The hard-cut rule for majority-water adjacencies (discussed in
    Section~\ref{sec:algorithm}) implements a well-established legal doctrine:
    that water bodies are natural community boundaries and that cross-water
    districts require specific justification.
  \item \textbf{Partisan neutrality.}
    No electoral, racial, or demographic data enters the computation.
    The signal is the physical character of the land surface.
\end{enumerate}

\paragraph{Distinction from M.3 housing character.}
The M.3 paper \citep{m3housing} operationalizes housing character through
ACS survey data on housing tenure, age, and value.
M.2 operationalizes the physical structure of what is built.
The signals overlap---residential development concentration (M.2)
correlates with single-family housing density (M.3)---but are not
redundant.
M.2 captures the spatial pattern of development types as seen from satellite;
M.3 captures the socioeconomic character of the households living in
that development.
Both contribute to the M.8 composite index \citep{m8composite}.

\paragraph{Paper organization.}
Section~\ref{sec:background} describes the NLCD 2021 data and the
physical geography of communities of interest.
Section~\ref{sec:algorithm} defines the land use character vector,
cosine similarity, hard-cut rule, and edge weight formula.
Section~\ref{sec:results} describes the experimental design and
defers quantitative results to Phase~2.
Section~\ref{sec:legal} provides a legal usage guide including the
water boundary doctrine and natural boundary case law.
Section~\ref{sec:conclusion} concludes.
